\documentclass[11pt]{article}

\usepackage[a4paper,margin=1in]{geometry}
\usepackage{amsmath,amssymb,amsthm,mathtools}
\usepackage{booktabs}
\usepackage{microtype}
\usepackage{enumitem}
\usepackage{hyperref}

\hypersetup{colorlinks=true,linkcolor=blue,citecolor=blue,urlcolor=blue}

\newtheorem{theorem}{Theorem}[section]
\newtheorem{proposition}[theorem]{Proposition}
\newtheorem{lemma}[theorem]{Lemma}
\newtheorem{corollary}[theorem]{Corollary}
\theoremstyle{remark}
\newtheorem{remark}[theorem]{Remark}
\newtheorem{conjecture}[theorem]{Conjecture}

\newcommand{\C}{\binom}
\newcommand{\R}{\mathbb{Q}}
\newcommand{\N}{N}
\newcommand{\X}{X}
\newcommand{\elln}{\ell}
\newcommand{\G}{G_{k,\ell}}
\newcommand{\Lone}{L_1}
\newcommand{\Ltwo}{L_2}
\newcommand{\Lthree}{L_3}

\title{\textbf{Exact Symmetric Homogeneous Layers of the \(p,q\)-Kernel}}\author{}
\date{18 August 2026}

\begin{document}
\maketitle

\begin{abstract}
A sequence of exact symbolic experiments was used to study a polynomial kernel
\(F_{k,\ell}(p,q)\) through the elementary symmetric variables
\(N=pq\) and \(S=p+q\), without assuming the parity-restricted factorization by
\(p+q+1\).  Writing
\(X=S+1=p+q+1\), the kernel admits a homogeneous expansion in \(N,X\).
The main structural result is an exact peeling hierarchy: the top homogeneous layer,
the complete first residual layer, and the complete second residual layer can be written
in closed form.  The third layer is known to be cubic in \(\ell\), with exact binomial laws
for its leading coefficient and several other components, while one coefficient component
remains unresolved in the current experimental record.  The original theorem-index
identification problem is therefore substantially reduced but not yet solved.
All claims labeled ``exact'' below are identities reconstructed and checked directly from
the symbolic \(p,q\) kernel over exact rational arithmetic; empirical interpolation is kept
separate from proved structure.
\end{abstract}

\section{Problem and notation}

Let \(F_{k,\ell}(p,q)\) denote the exact polynomial kernel under investigation, where
\(k\) is odd in the principal data set and \(\ell\) is the theorem parameter.  The
experiments showed that the first natural symmetric variables are
\[
    N=pq,\qquad S=p+q,
\]
and, for the homogeneous analysis, the shifted variable
\[
    X=S+1=p+q+1.
\]
The resulting symmetric polynomial is written
\[
    G_{k,\ell}(N,X)
    =F_{k,\ell}(p,q)\big|_{N=pq,\,X=p+q+1}.
\]
No division by \(p+q+1\) is used in this definition.

A recurring parity phenomenon was established first.  For the tested grid,
\[
 p+q+1 \mid F_{k,\ell}(p,q)
\quad\Longleftrightarrow\quad
 k\equiv \ell\pmod 2.
\]
For opposite parity the remainder is exactly the boundary specialization
\[
 F_{k,\ell}(-q-1,q),
\]
which is nonzero in all tested cases.  Thus the quotient
\(F/(p+q+1)\) cannot be the universal object.  This motivated the direct symmetric
reduction to \(G_{k,\ell}(N,X)\).

\section{Parity obstruction and direct symmetric reduction}

The direct reduction \(F\mapsto G\) was independently reconstructed on many examples,
including both parity branches.  The reconstruction audit had zero failures in the
reported experiments.  In particular, the parity obstruction is isolated in the
specialization \(X=0\), equivalently \(S=-1\), rather than in the symmetric reduction
itself.

Returning temporarily to \(S\), one may write
\[
    G_{k,\ell}(N,S)
    =(S+1)H_{k,\ell}(N,S)+R_{k,\ell}(N),
\]
where
\[
    R_{k,\ell}(N)=G_{k,\ell}(N,-1).
\]
For same parity, the experiments found \(R_{k,\ell}=0\); for opposite parity,
\(R_{k,\ell}\neq0\) and agrees with the exact obstruction obtained from
\(F(-q-1,q)\).  This gives a parity-unified decomposition while preserving the full
symmetric kernel.

\section{Homogeneous peeling in \(N,X\)}

The decisive change of viewpoint was to decompose \(G_{k,\ell}\) by total degree:
\[
    G_{k,\ell}=G^{(\ell)}+G^{(\ell-1)}+G^{(\ell-2)}+G^{(\ell-3)}+\cdots .
\]
We denote the first four layers by
\[
    G_{\mathrm{top}}=G^{(\ell)},
    \qquad
    \Lone=G^{(\ell-1)},
    \qquad
    \Ltwo=G^{(\ell-2)},
    \qquad
    \Lthree=G^{(\ell-3)}.
\]
This ordering is important: the layers are derived from the exact kernel first, and
only afterwards compared with theorem anchors.

\subsection{The top layer}

The corrected boundary audit established the coefficient identity
\[
    [N^aX^{\ell-a}]\,G_{\mathrm{top}}
    =-\C{k}{a},
    \qquad 0\le a\le k-1.
\]
Equivalently,
\begin{equation}
\boxed{
    G_{\mathrm{top}}
    =-X^{\ell-k}\bigl((X+N)^k-N^k\bigr).
}
\label{eq:top}
\end{equation}
The equality was checked across the reported odd \(k\) families and both same- and
opposite-parity values of \(\ell\).  This is a structural identity, not an
anchor-fitting ansatz.

\subsection{The first residual layer}

After subtracting \(G_{\mathrm{top}}\), the degree \(\ell-1\) layer has support
\(0\le a\le k\), with coefficients defined by
\[
    \Lone=\sum_{a=0}^{k} A^{(1)}_{k,\ell}(a)\,N^aX^{\ell-1-a}.
\]
The exact interior law found in the experiments is
\begin{equation}
    A^{(1)}_{k,\ell}(a)
    =\ell\C{k+1}{a}-k\C{k}{a-1},
    \qquad 0\le a<k,
\label{eq:L1-interior}
\end{equation}
with the boundary correction
\begin{equation}
    A^{(1)}_{k,\ell}(k)
    =\ell(k+1)-k^2+k.
\label{eq:L1-boundary}
\end{equation}
Here and below we use the convention \(\C{n}{r}=0\) outside the usual range when
that is convenient for interior formulas.

The important point is that the first layer already exhibits the recurring pattern:
interior binomial structure plus a finite boundary correction.

\section{Complete second layer}

Write
\[
    \Ltwo
    =\sum_{a=0}^{k+1}D^{(2)}_{k,\ell}(a)
       N^aX^{\ell-2-a}.
\]
The exact support audit gives
\[
    \operatorname{supp}_N(\Ltwo)=\{0,1,\dots,k+1\}.
\]

\subsection{Interior law \(0\le a<k\)}

The exact quadratic law in \(\ell\) is
\begin{equation}
\boxed{
\begin{aligned}
D^{(2)}_{k,\ell}(a)
={}&-\frac12\C{k+2}{a}\,\ell^2\\
&+\C{k+2}{a}
  \frac{2(k+1)a+k+2}{2(k+2)}\,\ell\\
&-\frac{k\,a(a+1)}{2(k+2)}\C{k+2}{a},
\end{aligned}}
\label{eq:L2-interior}
\end{equation}
for \(0\le a<k\).
The leading coefficient and second finite difference were independently audited:
\[
    [\ell^2]D^{(2)}_{k,\ell}(a)
    =-\frac12\C{k+2}{a},
\]
and
\[
    \Delta_{\ell,2}^2 D^{(2)}_{k,\ell}(a)
    =-4\C{k+2}{a}.
\]
These identities were verified across several odd values of \(k\) and multiple
values of \(\ell\).

\subsection{Boundary \(a=k\)}

At \(a=k\), the quadratic and linear coefficients agree with the interior formula,
but the constant term receives the exact correction
\[
    \boxed{\binom{k+2}{2}\frac{k}{k+1}?}
\]
The experimentally established correction is more simply written as
\begin{equation}
    \boxed{\frac{k(k+1)}2},
\label{eq:L2-correction}
\end{equation}
added to the constant term of the interior extrapolation.  Equivalently,
\[
    D^{(2)}_{k,\ell}(k)
    =D^{(2),\mathrm{int}}_{k,\ell}(k)+\frac{k(k+1)}2.
\]
The correction was checked for \(k=3,5,7,9,11,13\) and several values of \(\ell\).

\subsection{True endpoint \(a=k+1\)}

The final nonzero coefficient obeys the exact relation
\begin{equation}
    \boxed{
    D^{(2)}_{k,\ell}(k+1)
    =\frac{2}{k+1}D^{(2)}_{k,\ell}(k).
    }
\label{eq:L2-endpoint}
\end{equation}
This relation was checked on all reported endpoint families.  The support terminates
there, so
\[
    D^{(2)}_{k,\ell}(a)=0,
    \qquad a\ge k+2.
\]
Thus the second layer is completely determined by the interior formula, its single
boundary correction at \(a=k\), and the final endpoint relation \eqref{eq:L2-endpoint}.

\section{Third layer: exact cubic dependence and what is proved}

The third residual layer is defined by
\[
    \Lthree
    =G^{(\ell-3)}
    =\sum_a D^{(3)}_{k,\ell}(a)N^aX^{\ell-3-a}.
\]
The experiments established that for each fixed \((k,a)\), the coefficient is cubic
in \(\ell\).  Equivalently,
\[
    \Delta_{\ell}^3D^{(3)}_{k,\ell}(a)
\]
is constant in \(\ell\).  The leading cubic coefficient was proved to satisfy
\begin{equation}
\boxed{
    [\ell^3]D^{(3)}_{k,\ell}(a)
    =\frac16\C{k+3}{a}.
}
\label{eq:L3-A}
\end{equation}
The reported data also support exact formulas for the linear and constant-in-\(\ell\) components in a substantial interior range, but the current record does not
justify promoting the full cubic formula to a theorem: several component audits
showed systematic discrepancies in the proposed closed form for the middle
coefficient.

Thus it is useful to record the experimentally stable decomposition
\[
    D^{(3)}_{k,\ell}(a)
    =A_{k,a}\ell^3+B_{k,a}\ell^2+C_{k,a}\ell+D_{k,a},
\]
with the exact law \eqref{eq:L3-A} for \(A_{k,a}\), while treating the proposed
closed formulas for \(C_{k,a}\) and any full interior package as unresolved.

\subsection{Boundary behavior of \(L_3\)}

A robust boundary phenomenon was established for \(a=k\): when the coefficient is
compared to the interior extrapolation, the exact additive correction is
\begin{equation}
    \boxed{\binom{k+2}{3}}.
\label{eq:L3-k-correction}
\end{equation}
This correction is independent of \(\ell\) in all tested families.

The next row \(a=k+1\) is nonzero and has a cubic dependence on \(\ell\).  The current
experimental record gives explicit polynomials for several odd \(k\), for example
\[
\begin{aligned}
 k=3:&\quad \frac52\ell^3-\frac{65}{2}\ell^2+120\ell-105,\\
 k=5:&\quad \frac{14}{3}\ell^3-\frac{175}{2}\ell^2+\frac{3017}{6}\ell-840,\\
 k=7:&\quad \frac{15}{2}\ell^3-\frac{369}{2}\ell^2+1437\ell-3465,
\end{aligned}
\]
with analogous exact data for \(k=9,11\).  These observations strongly suggest
that the final nonzero row has a simple continuation of the cubic structure, but a
uniform derivation was not completed in the supplied experiments.

The support observed so far is finite:
\[
    D^{(3)}_{k,\ell}(a)=0
    \quad\text{for sufficiently large }a,
\]
with the nonzero tail terminating after the two boundary rows in the tested
families.  The exact general support theorem should be derived rather than inferred
from finite data.

\section{What the experiments rule out}

Several natural candidate coordinates were tested against the theorem-index anchors
and found not to be structurally stable.

\begin{enumerate}[leftmargin=2em]
    \item The universal quotient \(F/(p+q+1)\) is impossible because divisibility
          holds only on the same-parity branch.
    \item Direct coefficients of \(G(N,S)\) do not identify the theorem index \(s\)
          on the tested anchors.
    \item Shifted coefficients in \(X=S+1\) produce isolated matches but no stable
          universal coordinate rule.
    \item Effective indices such as \(r=s-(k-1)/2\) and boundary distances such as
          \(d=\ell-s\) collide across anchors with different theorem values.
    \item The outer diagonal is structurally simple, but its value alone is not the
          theorem coefficient.
\end{enumerate}

These failures are useful: they show that the original theorem coefficient should be
sought in the exact layered structure, not by unrestricted coordinate fitting.

\section{Relation to the original theorem-index problem}

The original target was to identify a natural coefficient coordinate corresponding to
the theorem index \(s\).  That problem remains open in the present record.
Nevertheless, the search space has changed substantially.  The exact kernel now has a
canonical hierarchy
\[
\boxed{
G_{k,\ell}
=G_{\mathrm{top}}+\Lone+\Ltwo+\Lthree+\cdots
}
\]
where the first three layers have progressively stronger closed structure.

The theorem anchors used in the experiments included, among others,
\[
\begin{array}{c|c|c}
(k,\ell,s)&\text{parity branch}&\text{theorem value}\\
\hline
(1,10,2)&\text{opposite}&27\\
(3,20,6)&\text{opposite}&935\\
(7,40,15)&\text{opposite}&-1797818\\
(3,7,3)&\text{same}&-3\\
(3,9,4)&\text{same}&9\\
(3,11,5)&\text{same}&-16\\
(5,11,5)&\text{same}&-5\\
(5,19,9)&\text{same}&-196\\
(9,21,11)&\text{same}&-84
\end{array}
\]
No proposed simple monomial coefficient coordinate passed the full anchor set.
This is consistent with the hypothesis that the theorem coefficient is a boundary,
linear-combination, or transformed-layer functional rather than a single raw
coefficient of \(G\).

\section{Methodological conclusions}

The computational program has converged to a reliable workflow.

\begin{enumerate}[leftmargin=2em]
    \item Construct the \(p,q\) kernel exactly.
    \item Verify symmetry and reconstruct it in \((N,S)\), then shift to \(X=S+1\).
    \item Decompose by total homogeneous degree.
    \item Derive each layer from the exact kernel before inspecting theorem anchors.
    \item Treat boundary rows separately whenever the interior binomial law stops.
    \item Use exact arithmetic throughout; in particular, rational constants must be
          represented by exact rational objects rather than Python float division.
\end{enumerate}

This workflow fixed several recurring computational errors encountered during the
experiments, including invalid polynomial construction after symbolic division,
misuse of \texttt{NegativeInfinity}, and floating-point equality failures.

\section{Open problems}

The main mathematical questions left by the current record are:

\begin{enumerate}[leftmargin=2em]
    \item derive the complete closed form of \(L_3\), especially the unresolved
          middle coefficient(s), directly from the \(p,q\) kernel;
    \item prove the general boundary support and endpoint laws for all higher layers;
    \item determine whether every homogeneous layer has polynomial dependence on
          \(\ell\) of degree equal to its layer number;
    \item identify the exact theorem-index functional that turns the layered
          coefficients into the published theorem values;
    \item derive that functional symbolically rather than fitting it to the nine
          anchor values.
\end{enumerate}

\section{Conclusion}

The central outcome is a shift from coefficient hunting to exact structure.
The polynomial kernel is symmetric in \(p,q\), admits a parity-unified formulation in
\(N=pq\) and \(X=p+q+1\), and has an exact homogeneous peeling hierarchy.  The top
layer is binomial, the first residual layer is binomial-linear, the second residual
layer is binomial-quadratic with controlled boundary corrections, and the third
residual layer is cubic in \(\ell\) with an exact binomial leading coefficient.

The theorem-index problem has therefore become a much more constrained algebraic
question.  The next successful step should come from deriving the remaining layer
and then locating the theorem coefficient as a natural functional of this exact
hierarchy.

\bigskip
\noindent\textbf{Status summary.}
Top layer: exact.\\
First layer: exact with boundary correction.\\
Second layer: exact, including endpoints.\\
Third layer: exact cubic dependence and leading law established; full closed middle
coefficient still unresolved.\\
Theorem-index identification: unresolved.

\end{document}
